\section{Objetivos y KPIs}

\textbf{Objetivo general.} Analizar relaciones entre atributos musicales y contextuales y la
popularidad observada, dejando una base confiable y evidencia reproducible para evaluar modelos de
regresión.

\textbf{Objetivos específicos.}

\begin{enumerate}
\item comprender el problema, los interesados y las decisiones relevantes;
\item documentar procedencia, estructura y calidad de la fuente;
\item preparar una unidad de análisis de una canción por \code{track_id};
\item explorar distribuciones, relaciones, anomalías, redundancias y sesgos;
\item clasificar variables según utilidad y riesgo para modelamiento;
\item comunicar limitaciones éticas, de privacidad y de interpretación.
\end{enumerate}

\begin{table}[H]
\centering
\caption{Indicadores del proyecto y su función.}
\label{tab:kpis}
\begin{tabularx}{\textwidth}{p{2.6cm}p{4.2cm}X}
\toprule
Tipo & Indicador & Interpretación y uso \\
\midrule
Negocio & Media y mediana de popularidad & Describen el centro del catálogo observado; no son
metas de predicción por sí mismas. \\
Negocio & Proporción con popularidad $\geq 70$ & Segmento operativo de alta popularidad; el umbral
no es una categoría oficial de Spotify. \\
Negocio & Popularidad por género & Compara segmentos con soporte; no estima prevalencia real ni
efecto causal. \\
Calidad & Completitud & Controla celdas disponibles y evita imputaciones injustificadas. \\
Calidad & Duplicidad de \code{track_id} & Verifica unidad de análisis y previene fuga entre
entrenamiento y prueba. \\
Modelo & MAE & Error absoluto promedio en puntos de popularidad, interpretable en la escala original. \\
Modelo & RMSE & Penaliza más los errores grandes y complementa MAE. \\
Modelo & $R^2$ & Proporción de variabilidad explicada respecto de una línea base; no mide causalidad. \\
\bottomrule
\end{tabularx}
\end{table}

Los KPIs de negocio describen utilidad potencial; las estadísticas descriptivas caracterizan la
muestra; y las métricas técnicas solo serán válidas al evaluarse fuera de muestra.
